
In this chapter we will show that in the case of (well-behaved: locally Noetherian, ...) affine schemes the desired functors $f_!$ and $f^!$ exist and that they satisfy some expected properties. 

It is to note that in \cite{condenseddotpdf} Clausen and Scholze show this only in case of finite type affine $\Z$-schemes, \confer \cite[theorem 8.13]{condenseddotpdf}. While they later \emph{globalize} the result in \cite[Chapters X \& XI]{condenseddotpdf} to arbitrary schemes and hence recover the unrestricted affine case, this is somewhat unsatisfactory to the author. One then essentially describes an affine scheme $\Spec A$ as the limit $\limit_{A'} \Spec A'$ consisting of sequences of compatible points from the affine schemes $\Spec A'$ for finitely generated $\Z$-subalgebras of $A$. This is of course quite inexplicit and, to the author, not so geometrically appealing. Even though we will of course adopt the same viewpoint, we will do some more implicitly and try to develop the results given by Clausen and Scholze for general affine schemes instead. This entails describing many of the results given by Clausen and Scholze -- but after some amount of the above gluing. Proving the generalized results often worked, but not always! If there is some complication, we will denote it (obviously) at each place.

While \quote{explicitly} describing this gluing is of course technically redundant after one has already globalized, it is still somewhat more satisfying to the author.

Additionally, the mentioned globalization makes use of $\infty$-theoretical results, as one needs to glue derived categories of the form $\derived{(A', R')_\blacksquare}$ -- something that is not possible on the level of \quote{ordinary} derived categories. While we will state the results in chapter \ref{chapter: globalization}, we will not make an attempt to formally justify them. Hence, if we were to not translate the results of \cite[theorem 8.13]{condenseddotpdf} to general affine schemes, we would have, in this thesis, only proven duality for finite type affine $\Z$-schemes -- a very restrictive geometric setting (this excludes for example all affine schemes over infinite fields).

Before we start with our preparations we should fix some notation. Much like the category of affine schemes can be thought of as $\commutativeRings^\op$, it will prove useful to talk about the analytic rings $A_\blacksquare$ and $(A, R)_\blacksquare$ as if they were geometric objects. For this we will introduce some terminology. Note that this terminology is not used by Clausen \& Scholze.
\begin{definition}[Analytic spectra \& virtual maps]
	\label{definition: analytic spectra and virtual maps}
	
	Let $\Spec\colon \analyticRings \to \analyticRings^\op$ be the canonical contravariant functor. For an analytic ring $\preAnaRing{A}$ the object $\Spec \preAnaRing{A} \in \analyticRings^\op$ is called the \emph{analytic spectrum of $\preAnaRing{A}$}. Any map of analytic rings $\preAnaRing{A}\to \preAnaRing{B}$ induces a map of analytic spectra $\Spec\preAnaRing{B}\to \Spec\preAnaRing{A}$. Such a map of analytic spectra $\Spec \preAnaRing{B} \to \Spec \preAnaRing{A}$ is called \emph{virtual} if the associated morphism $\underlying{A}\to \underlying{B}$ is the identity. 
\end{definition}

\begin{remark}
	If one is acquainted with the theory of adic spaces this might seem familiar (see chapter \ref{chapter: globalization} for a very brief introduction). Indeed, to any pair of rings $R\subseteq A$ with $R$ integrally closed in $A$ one can associate the adic spectrum $\Spa(A, R)$. Furthermore, if $R\to A$ is general, and we denote by $\tilde R$ the integral closure of (the image of) $R$ inside $A$, then by lemma \ref{lemma: analytic rings associated to finite morphisms} we find that $R_\blacksquare \cong \tilde R_\blacksquare$ and hence that $(A, R)_\blacksquare \cong (A, \tilde R)_\blacksquare$. That is for a general $(A, R)_\blacksquare$ one can associate the adic spectrum $\Spa(A, \tilde R)$. This seems to be a proper dictionary, but the author has not investigated it. As we will not make use of the theory of adic space until chapter \ref{chapter: globalization} it is simply more convenient to introduce artificial language as in definition \ref{definition: analytic spectra and virtual maps}.
\end{remark}

\begin{remark}[Notation]
	Suppose that $\phi\colon \preAnaRing{A}\to \preAnaRing{B}$ is a morphism of analytic rings with associated map of analytic spectra $f\colon \Spec \preAnaRing{B} \to \Spec\preAnaRing{A}$. To respect the change in perspective when switching from analytic rings (the algebraic objects) to analytic spectra (the geometric objects) we will denote by $f_\ast$ the scalar restriction $\scalarRestriction{\phi}\colon \mod{\preAnaRing{B}}\to\mod{\preAnaRing{A}}$ and by $f^\ast$ the scalar extension $\completedScalarExtension{\phi}{\preAnaRing{B}}\colon \mod{\preAnaRing{A}}\to\mod{\preAnaRing{B}}$. Since we will be mostly working in the derived categories we will further write $f_\ast$ and $f^\ast$ for the associated derived functors. 
\end{remark}

\begin{remark}[What about virtual maps?]
	Any morphism of discrete rings $\phi\colon R\to A$ (corresponding to morphism of schemes $f\colon \Spec A\to \Spec A$) gives rise to three (in general distinct) analytic rings: $R_\blacksquare$, $(A, R)_\blacksquare$ and $A_\blacksquare$. In fact, these provide a factorization of the map of analytic rings $\phi\colon R_\blacksquare \to A_\blacksquare$ as
		$$R_\blacksquare \xrightarrow{\phi} (A, R)_\blacksquare \xrightarrow{\id_A} A_\blacksquare.$$
	If one now wants to study how modules over $R$ and $A$ (\ie quasi-coherent sheaves on $\Spec R$ and $\Spec A$) behave under transport along $\phi$ the main difficulty comes from the second map $\id_A\colon (A, R)_\blacksquare\to A_\blacksquare$ -- with associated virtual map $j\colon \Spec A_\blacksquare\to\Spec(A, R)_\blacksquare$. Indeed, $(A, R)_\blacksquare$ is essentially \emphquote{just a copy of $R_\blacksquare$ in the world of $A$-modules}. However, the difference between $(A, R)_\blacksquare$ and $A_\blacksquare$ is surprisingly subtle and most of the remaining work until proving the existence of a good theory of $f_!$ and $f^!$ in theorem \ref{theorem: exceptional direct and inverse image functors, the affine case} is dedicated to the study of this virtual map in the special case of $A=R[T]$, the coordinate ring of the affine space $\A_R^1$ over $\Spec R$. For example, we will construct a functor $j_!$ that will be used to define an appropriate exceptional direct image $f_!\colon \derived{A_\blacksquare}\to\derived{R_\blacksquare}$ that gives rise to $f^!$ as its right adjoint.
\end{remark}

\section{Preparations: The case of \texorpdfstring{$\A_R^1$}{Affine Space over an Affine Base}}

As already seen in the proof of corollary \ref{corollary: analyticity of Ablack for finitely generated free Z-algebras} and theorem \ref{theorem: analyticity of finite type Z-algebras}, it is often relatively straightforward to reduce the proof a given claim from the general case of finite type algebra to a free algebra to the free algebra in one variable. We will do the same in for example the proof of theorem \ref{theorem: exceptional direct and inverse image functors, the affine case} which asserts the existence of $f_!$ and $f^!$. We will thus study this case thoroughly. To this end consider the following convention.
\begin{convention}
	Throughout this section let $R$ be a discrete ring. Denote by $A\ldef R[T]$ the ring of polynomials. Observe that the map of rings $R_\blacksquare \to A_\blacksquare$ factors as
		$$R_\blacksquare \to (A, R)_\blacksquare \to A_\blacksquare$$
	and denote by $j\colon \Spec A_\blacksquare \to \Spec (A, R)_\blacksquare$ the virtual map.
%	with $j_\ast\colon \derived{A_\blacksquare} \to \derived{(A, R)_\blacksquare}$ $j^\ast = A_\blacksquare\Dotimes_{(A, R)_\blacksquare} -$ is associated scalar restriction and scalar extension respectively.
\end{convention}

We will now translate many of the result from section \ref{section: an analogue of hilbert's basis theorem} to this more general case where $R$ is an arbitrary discrete ring and no longer $\Z$-finitely generated.

It turns out that the analogue of $A_\infty$ is not as straightforward. One has to glue.
\begin{definition}
	Denote by
	$$A_\infty\ldef \colimit_{R'\subseteq R} \laurent{R'}{T^{-1}}$$
	the subset of $\laurent{R}{T^{-1}}$ of formal Laurent series $\sum_{n\gg \infty} r_n T^{-n}$ in $T^{-1}$ all of whose coefficients $(r_n)_{n\in \Z}$ lie in some finitely generated $\Z$-algebra $R'$ of $R$.
	
	Then $A_\infty$ is a topological $A$-subalgebra of $\laurent{R}{T^{-1}}$ with the natural $\ideal{T^{-1}}$-adic topology. Indeed, $A_\infty$ is a ring, since if $x \in \laurent{R'}{T^{-1}}$ and $y \in \laurent{R''}{T^{-1}}$ then $x\cdot y\in \laurent{\Z[R'\cup R'']}{T^{-1}}\subseteq A_\infty$ where $\Z[R'\cup R']$ denotes the $\Z$-subalgebra of $R$ generated by $R'$ and $R''$. Furthermore, since each $f\in A=R[T]$ has only finitely many coefficients it is clear that $f\in A_\infty$. Hence $A\subseteq A_\infty$ and $A_\infty$ is an $A$-algebra. Even more, $A_\infty$ is then an $A$-subalgebra of the topological $A$-algebra $\laurent{R}{T^{-1}}$ and thus inherits the $\ideal{T^{-1}}$-adic topology.
\end{definition}

\begin{remark}
	If $R$ is a finitely generated $\Z$-algebra then the definition of $A_\infty$ drastically simplifies. Indeed, for every formal Laurent series over $R$, its coefficient lie in $R$, which is finitely generated. Hence, $A_\infty$ is simply $\laurent{R}{T^{-1}}$.
\end{remark}

There is a useful way to rewrite $A_\infty$.
\begin{lemma}[Alternative representation of $A_\infty$]
	\label{lemma: alternative representation of Ainfty}
	
	The morphism
		$$A_\infty = \colimit_{R'\subseteq R} \laurent{R'}{T^{-1}} \to \colimit_{R'\subseteq R} A\otimes_{R'[T]} \laurent{R'}{T^{-1}}$$
	given by mapping $s \in \laurent{R'}{T^{-1}}$ to $1\otimes s\in A\otimes_{R'[T]} \laurent{R'}{T^{-1}}$ is an $A$-module isomorphism.
\end{lemma}
\begin{proof}
	Consider the morphism given by mapping $f\otimes s$ with $f = \sum_{m=0}^M r_m T^m \in A=R[T]$ and $s=\sum_{n\gg -\infty} r_nT^{-n} \in \laurent{R'}{T^{-1}}$ to $f\cdot s \in \laurent{R}{T^{-1}}$. If its image is contained in $A_\infty$ then it is clear that multiplication is the desired inverse. Indeed, both $f$, $s$ and hence $f\cdot s$ are contained in the subalgebra
		$$\laurent{\Z[\{r_0, \dots, r_M\}\cup R']}{T^{-1}}$$
	of $A_\infty$ and the claim follows.
	
	Via a similar argument one shows that the given (clearly additive) morphism is $A$-linear. Indeed, for $f=\sum_{m=0}^M r_mT^m\in A=R[T]$ any polynomial and $s\in\laurent{R'}{T^{-1}}$ any series, the element $f\cdot s \in A_\infty$ can be considered in $\laurent{\Z[\{r_0,\dots, r_M\}\cup R'\}]}{T^{-1}}\subseteq A_\infty$ instead and then $1\otimes f\cdot s = f\cdot s = f\cdot(1\otimes s)$ as the tensor product has the right balance. Its inverse (as a map of sets) is then automatically $A$-linear as well.
\end{proof}

\begin{lemma}[Flatness]
	The $A$-module $A_\infty$ is a $A$-flat.
\end{lemma}
\begin{proof}
%	Observe that the natural morphism of $A$-modules
%	$$A_\infty = \colimit_{R'\subseteq R} \laurent{R'}{T^{-1}}\isorightarrow \colimit_{R'\subseteq R} R\otimes_{R'} \laurent{R'}{T^{-1}}$$
%	is an isomorphism as it admits an inverse by mapping $r\otimes \sum_n r'_n T^{-n}$ in $R\otimes_{R'}\laurent{R'}{T^{-1}}$ to $\sum_n rr'_n T^{-n}$ in $\laurent{\Z[R', r]}{T^{-1}}$. 
	
	By lemma \ref{lemma: alternative representation of Ainfty} we know that
		$$A_\infty \cong \colimit_{R'\subseteq R} R[T]\otimes_{R'[T]} \laurent{R'}{T^{-1}}$$
	as $A$-modules. Observe that each $R[T]\otimes_{R'[T]} \laurent{R'}{T^{-1}}$ is the base change of the $R'[T]$-module $\laurent{R'}{T^{-1}}$ along the inclusion $R'[T]\subseteq R[T]$. But the $R'[T]$-module $\laurent{R'}{T^{-1}}$ is flat by lemma \ref{lemma: Ainfty is A-flat, f.g. case} and then so is its base change by \cite[\href{https://math.stackexchange.com/a/120839/529214}{this argument}]{zhenlinbasechange} of \href{https://math.stackexchange.com/users/5191/zhen-lin}{Zhen Lin}. Hence, $A_\infty$ is a filtered colimit of flat $A$-modules and hence flat as a consequence of Lazard's theorem \cite[\href{https://stacks.math.columbia.edu/tag/058G}{Tag 058G}]{stacks-project}.
\end{proof}

\begin{lemma}
	\label{lemma: presentation for Afinty, general case}
	
	The topological $A$-module $A_\infty$ fits into the short exact sequence
		$$0\to \colimit_{R'\subseteq R} \series{R'}{U}[T]\xrightarrow{U\cdot T - 1}\colimit_{R'\subseteq R} \series{R'}{U}[T] \to A_\infty \to 0$$
	of $A$-modules where $\colimit_{R'\subseteq R} \series{R'}{U}[T]$ is a compact projective object of $\mod{(A, R)_\blacksquare}$. Hence, is $(A, R)_\blacksquare$-complete, $R_\blacksquare$-complete and quasi-isomorphic to the bounded complex of compact projectives of $\mod{(A, R)_\blacksquare}$ so is a derived compact in $\derived{(A, R)_\blacksquare}$.
\end{lemma}
\begin{proof}
	Observe that for any $R'$ there is a short exact sequence
	$$0\to \series{R'}{U}[T]\xrightarrow{UT-1} \series{R'}{U}[T] \to \laurent{R'}{T^{-1}}\to 0$$
	of $R'$-modules since $\laurent{R'}{T^{-1}} \cong \series{R'}{T^{-1}}[T]$. These induce an exact sequence
	$$0\to\colimit_{R'\subseteq R} \series{R'}{U}[T] \xrightarrow{UT-1} \colimit_{R'\subseteq R} \series{R'}{U}[T] \to \underbrace{\colimit_{R'\subseteq R} \laurent{R'}{T^{-1}}}_{=A_\infty}\to 0$$
	of condensed abelian groups, since filtered colimits in $\condensedAb$ are exact. As each of these groups is a condensed $A$-module and as each of the given morphisms is $A$-linear this is even an exact sequence of condensed $A$-modules. Analogous to lemma \ref{lemma: alternative representation of Ainfty} we see that
		$$\colimit_{R'\subseteq R} \series{R'}{U}[T] \cong \colimit_{R'\subseteq R} A\otimes_{R'[T]}\series{R'}{U}[T]\cong \colimit_{R'\subseteq R} A\otimes_{R'}\series{R'}{U}$$
	as $A\otimes_{R'[T]} \series{R'}{U}[T]\cong A\otimes_{R'[T]} (R'[T]\otimes_{R'}\series{R'}{U}) \cong A\otimes_{R'}\series{R'}{U}$. But $\series{R'}{U} \cong \prod_\N R'$ and hence it follows by (the proof of) lemma \ref{lemma: more compact projectives} that
		$$\colimit_{R'\subseteq R} \series{R'}{U}[T] \cong A\otimes_{R}\colimit_{R'\subseteq R} R\otimes_{R'} \prod_\N R'$$
	is a compact projective object of $\mod{(A, R)_\blacksquare}$ (as the retraction is still valid after application of $A\otimes_{R} -$). Then the claim follows at once.
	
%	Denote by $\hat \N$ the one-point-compactification of $\N$ and observe that the sequence of $R'$-modules
%	$$0\to \free{R'_\blacksquare}{\infty}\to\free{R'_\blacksquare}{\hat \N} \to \series{R'}{U} \to 0$$
%	where the first morphism is induced by the inclusion $\{\infty\}\to \hat\N$ and where the second morphism corresponds to the map $\hat \N\to\series{R'}{U}$ given by $n\mapsto U^n$ and $\infty\mapsto 0$, is exact. Denote the inclusion $\free{R'_\blacksquare}{\infty} \to \free{R'_\blacksquare}{\hat\N}$ by $s$ and denote by $r\colon \free{R'_\blacksquare}{\hat\N}\to \free{R'_\blacksquare}{\infty}$ the morphism induced by the projection $\hat\N\to\{\infty\}$. By functoriality of $\free{R'_\blacksquare}{-}$ we obtain that $r\circ s = \id$ as $\{\infty\}\to \hat\N\to\{\infty\}$ is the identity. In particular $r$ is a retraction for $s$ and the short exact sequence is split. As any additive functor preserves split exact sequences, the induced sequence
%	$$0\to \underbrace{\colimit_{R'\subseteq R} A\otimes_{R'} \free{R'_\blacksquare}{\infty}}_{\cong\free{(A, R)_\blacksquare}{\infty}}\to \underbrace{\colimit_{R'\subseteq R} A\otimes_{R'} \free{R'_\blacksquare}{\hat\N}}_{\cong\free{(A, R)_\blacksquare}{\hat\N}}\to \colimit_{R'\subseteq R} A\otimes_{R'} \series{R'}{U}\to 0$$
%	is split exact as well, implying that $\colimit_{R'\subseteq R} A\otimes_{R'}\series{R'}{U}$ is a retract of $\free{(A, R)_\blacksquare}{\hat\N}$. But the latter is compact projective in $\mod{(A, R)_\blacksquare}$ and so must be its retract. 
	
%	As for any complex $C^\bullet\in\derived{(A, R)_\blacksquare}$ the functor $\Hom(-, C^\bullet)\colon \derived{(A, R)_\blacksquare}\to \Ab$ is cohomological by \cite[\href{https://stacks.math.columbia.edu/tag/0149}{Tag 0149}]{stacks-project} and the first two terms of the short exact sequence are compact, an application of the $5$-lemma yields, that $\colimit_{R'\subseteq R} A\otimes_R\series{R}{U}$ is compact as well.
%	
%	$$\begin{tikzcd}
%		&0\ar[d] &0\ar[d] &0\ar[d]\\
%		0\ar[r] & \free{(A, R)_\blacksquare}{\infty}\ar[r]\ar[d] &\free{(A, R)_\blacksquare}{\infty}\ar[r]\ar[d] &0\ar[d]\ar[r] &0\\
%		0\ar[r] &\free{(A, R)_\blacksquare}{\hat\N}\ar[r, "1_\N\cdot T - 1"]\ar[d, "n\mapsto U^n"] &\free{(A, R)_\blacksquare}{\hat\N}\ar[r]\ar[d, "n\mapsto U^n"]&Q\ar[d]\ar[r] &0\\
%		0\ar[r] &\colimit_{R'\subseteq R} \series{R'}{U}[T] \ar[r, "UT-1"]\ar[d] &\colimit_{R'\subseteq R} \series{R'}{U}[T]\ar[r]\ar[d] &A_\infty\ar[r]\ar[d] &0\\
%		&0 &0 &0
%	\end{tikzcd}$$
\end{proof}

\begin{lemma}[$A_\infty$ is idempotent]
	
	The $A$-algebra $A_\infty$ is idempotent, by which we mean that the multiplication map of $A_\infty$ gives rise to an isomorphism $A_\infty\Dotimes_{(A, R)_\blacksquare} A_\infty \to A_\infty$ in $\derived{(A, R)_\blacksquare}$.
\end{lemma}
\begin{proof}
	Recall that by lemma \ref{lemma: formal Laurent series are idempotent} we know that multiplication $\laurent{R}{T^{-1}}\otimes_{R[T]}\laurent{R}{T^{-1}}\to \laurent{R}{T^{-1}}$ is an $A$-module isomorphism. As both multiplication $a\otimes b \mapsto a\cdot b$ and its inverse $a\mapsto 1\otimes b$ restrict to morphisms between $A_\infty\otimes_A A_\infty$ and $A_\infty$ we immediately obtain that multiplication $A_\infty\otimes_A A_\infty\isorightarrow A_\infty$ for $A_\infty$ is an isomorphism as well.
	
	It remains to show that the derived multiplication is an isomorphism. This follows verbatim as in lemma \ref{lemma: Ainfty is idempotent, f.g. case} -- the case where $R$ was $\Z$-finitely generated.
\end{proof}

We again obtain by idempotentcy that $A_\infty$-modules are well-behaved.
\begin{lemma}[$A_\infty$-modules]
	\label{lemma: Ainfty-modules are reflective}
	
	The $A_\infty$-modules in $\derived{(A, R)_\blacksquare}$ form a full subcategory whose inclusion has a left adjoint given by the functor $A_\infty\Dotimes_{(A, R)_\blacksquare} -$. Furthermore, any $M \in \derived{(A, R)_\blacksquare}$ admits at most one $A_\infty$-module structure and such a structure exists if and only if the natural morphism $M\to A_\infty\Dotimes_{(A, R)_\blacksquare} M$ is an isomorphism.
\end{lemma}
\begin{shortproof}
	The proof of lemma \ref{lemma: Ainfty-modules, f.g. case} applies verbatim.
\end{shortproof}

\begin{lemma}[$A_\infty$-modules admit not extensions by $A_\blacksquare$-complete modules]
	\label{lemma: Ainfty-modules admit not extensions by Ablack-completes}
	
	For any $S$ extremally disconnected we have
		$$\eRHom_A(A_\infty, \free{A_\blacksquare}{S}) \cong 0$$
	in $\derived{\condensedAb}$. In particular, if $C^\bullet \in \derived{A_\blacksquare}$ and $M$ is any $A_\infty$-module, then $\eRHom_A(M, C^\bullet)\cong 0$.
\end{lemma}
\begin{proof}
	Recall that
	$$A_\infty = \colimit_{R'\subseteq R} \laurent{R'}{T^{-1}}\cong \colimit_{R'\subseteq R} R[T] \otimes_{R'[T]} \laurent{R'}{T^{-1}} = \colimit_{R'\subseteq R} A \otimes_{R'[T]} \laurent{R'}{T^{-1}}.$$
	As each $\laurent{R'}{T^{-1}}$ is $R'[T]$-flat by lemma \ref{lemma: Ainfty is A-flat, f.g. case} we find further that
		$$A_\infty \cong \colimit_{R'\subseteq R} A\Dotimes_{R'[T]} \laurent{R'}{T^{-1}}.$$
	Thus,
	\begin{align*}
		&\, \eRHom_A(A_\infty, \free{A_\blacksquare}{S})\\
		\cong &\, \eRHom_A(\colimit_{R'\subseteq R} A\Dotimes_{R'[T]} \laurent{R}{T^{-1}}, \free{A_\blacksquare}{S})\\
		\cong &\, \homotopyLimit_{R'\subseteq R} \eRHom_A(A\Dotimes_{R'[T]} \laurent{R}{T^{-1}}, \free{A_\blacksquare}{S})\\
		\cong &\, \homotopyLimit_{R'\subseteq R} \eRHom_{R'[T]}(\laurent{R}{T^{-1}}, \free{A_\blacksquare}{S})\\
	\end{align*}
	as $\eRHom_A$ preserves (homotopy) colimits by \assumptionPreWord \ref{assumption: right derived Hom maps colimits to limits} and as the $\condensedAb$-enriched adjunction of scalar extension and restriction along $R'[T]\to R[T]=A$ is derivable by lemma \ref{lemma: enriched adjunctions are stable under derivation}. Fix some $R'\subseteq R$. Recall that $\free{A_\blacksquare}{S}$ is defined as a colimit over the poset of all finitely generated $\Z$-subalgebras of $A=R[T]$. Analogous to the argument in the proof of theorem \ref{theorem: analyticity of Ablack for non-finite Z-algebras A} we observe that the sub-poset given by rings $R''[T]$ with $R'\subseteq R''\subseteq R$ some intermediate ring is cofinal. Then
	$$\free{A_\blacksquare}{S} \cong \colimit_{R'\subseteq R''\subseteq R} \free{(A, R''[T])_\blacksquare}{S}.$$
	But since $R''[T]_\blacksquare$ is analytic we obtain by lemma \ref{lemma: relative completeness} that each $\free{(A, R''[T])_\blacksquare}{S}$ is $R''_\blacksquare$-complete and hence also $R'[T]_\blacksquare$-complete by restricting along the map of analytic rings $R'[T]\to R''[T]_\blacksquare$ given by the inclusion. Thus, $\free{A_\blacksquare}{S}$ is a colimit of $R'[T]_\blacksquare$-complete modules and hence $R'[T]_\blacksquare$-complete itself. We thus obtain that for any $R'$ already
	$$\eRHom_{R'[T]}(\laurent{R'}{T^{-1}}, \free{A_\blacksquare}{S}) \cong 0$$
	by lemma \ref{lemma: Ainfty-modules admit no extensions, f.g. case}. It immediately follows that
	$$\eRHom_A(A_\infty, \free{A_\blacksquare}{S}) \cong 0$$
	as well.
	
	The remaining proof now consists of extending the result to any $A_\infty$-module $M$ and any $A_\blacksquare$-complete $C^\bullet$. This is analogous to the corresponding part in the proof of the finitely generated case \ref{lemma: Ainfty-modules admit no extensions, f.g. case}.
	
%	\begin{todo}
%		Does compactness imply commutation also for enrichments?
%	
%		Need the $A_\infty$-module lemma for the general $R$ case
%	\end{todo}
%		
%	Now if $C^\bullet = C[0]$ is concentrated in degree $0$ where $C = \bigoplus_{i\in I} \free{A_\blacksquare}{S_i}$ is $A_\blacksquare$-free then
%		$$\eRHom_A(A_\infty, C^\bullet) \cong \bigoplus_{i\in I} \eRHom_A(A_\infty, \free{A_\blacksquare}{S_i})\cong 0$$
%	as $A_\infty$ is derived compact in $\derived{(A, R)_\blacksquare}$ and as all $A_\blacksquare$-complete modules are $(A, R)_\blacksquare$-complete. By induction it follows that any bounded complex $C^\bullet$ of $A_\blacksquare$-free modules also satisfies $\eRHom_A(A_\infty, C^\bullet) \cong 0$. Now if $C^\bullet\in\derived{A_\blacksquare}$ is bounded to the right then there is a left resolution $P^\bullet \to C^\bullet$ by $A_\blacksquare$-free modules which is also bound to the right. Writing $P^\bullet$ as a sequential colimit of its truncations $P_n^\bullet \ldef \sTruncation{\ge -n} P^\bullet$ we obtain that
%		$$\eRHom_A(A_\infty, C^\bullet) \cong \eRHom_A(A_\infty, P^\bullet) \cong \colimit_{n\in \N} \eRHom_A(A_\infty, P_n^\bullet) \cong 0$$
%	by lemma \ref{lemma: commutation with sequential colimits}. Finally if $C^\bullet \in \derived{A_\blacksquare}$ is general, we can write it as the limit $\colimit_{n\in \N} \sTruncation{\le n}C^\bullet$ of bound to the right complexes which implies
%		$$\eRHom_A(A_\infty, C^\bullet)\cong \colimit_{n\in \N} \eRHom_A(A_\infty, \sTruncation{\le n} C^\bullet) \cong 0$$
%	as $\eRHom_A(A_\infty, -)$ preserves (homotopy) colimits by lemma \ref{lemma: commutation with sequential colimits} as well.
%	
%	
%	\begin{todo}
%		The following argument doesnt work for as we are only enriched over $\Ab$.
%	\end{todo}
%	Now suppose that $M$ is an $A_\infty$-module. Then $M\cong M\Dotimes_{(A, R)_\blacksquare} A_\infty$ by lemma \ref{lemma: Ainfty-modules, f.g. case}. Thus for any $C^\bullet \in \derived{A_\blacksquare}$ we find that
%	\begin{align*}
%		&\,\eRHom_A(M, C^\bullet) \\ \cong&\, \eRHom_A(M\Dotimes_{(A, R)_\blacksquare} A_\infty, C^\bullet)\\
%		\cong &\,\eRHom_A(\derivedCompletion{M\Dotimes_A A_\infty}{(A, R)_\blacksquare}, C^\bullet)\\
%		\cong &\,\eRHom_A(M\Dotimes_A A_\infty, C^\bullet)\\
%		\cong &\,\eRHom_A(M, \intRHom_A(A_\infty, C^\bullet)).
%	\end{align*}
%	But as previously shown $\intRHom_A(A_\infty, C^\bullet) \cong 0$ as the underlying (complex of) condensed abelian group(s) of $\intRHom_A$ is $\eRHom_A$. Thus
%	$$\eRHom_A(M,C^\bullet) \cong \eRHom_A(M, 0) \cong 0.$$
	
\end{proof}

\begin{lemma}
	\label{lemma: special cokernels are Ainfty-modules, general case}
	
	For any extremally disconnected $S$, the cokernel of the morphism
		$$\free{(A, R)_\blacksquare}{S} \to \free{A_\blacksquare}{S}$$
	is an $A_\infty$-module.
\end{lemma}
\begin{proof}
	
	Denote the cokernel by $Q$ and recall that $\free{R_\blacksquare}{S} = \colimit_{R'\subseteq R} R\otimes_{R'} \free{R'_\blacksquare}{S}$. Then
	\begin{align*}
		\free{(A, R)_\blacksquare}{S} &= A\otimes_{R} \colimit_{R'\subseteq R} R\otimes_{R'}\free{R'_\blacksquare}{S}\\
		&\cong  \colimit_{R'\subseteq R} A\otimes_{R'}\free{R'_\blacksquare}{S}\cong \colimit_{R'\subseteq R} A\otimes_{R'[T]} R'[T] \otimes_{R'} \free{R'_\blacksquare}{S}.
	\end{align*}
	 Now since the poset of polynomial rings $R'[T]$ with $R'\subseteq R$ a finitely generated $\Z$-subalgebra is cofinal in the poset of finitely generated $\Z$-subalgebras $A'\subseteq A$ we obtain that $\free{A_\blacksquare}{S} \cong \colimit_{R'\subseteq R} A\otimes_{R'[T]} \free{R'[T]_\blacksquare}{S}$. We now have for every $R'$ a short exact sequence
	 	$$0\to R'[T]\otimes_{R'}\free{R'_\blacksquare}{S} \to \free{R'[T]_\blacksquare}{S} \to Q_{R'} \to 0$$
	 of $R'[T]$-modules, where the cokernel $Q_{R'}$ is a $\laurent{R'}{T^{-1}}$-module by lemma \ref{lemma: special cokernels are Ainfty-modules, f.g. case}. As tensor products are right exact we obtain that the induced sequences
	 	$$A\otimes_{R'[T]} R'[T]\otimes_{R'} \free{R'_\blacksquare}{S} \to A\otimes_{R'[T]} \free{R'[T]_\blacksquare}{S} \to A\otimes_{R'[T]} Q_{R'} \to 0$$
	 of $A$-modules are exact as well and that each $A\otimes_{R'[T]} Q_{R'}$ is still an $\laurent{R'}{T^{-1}}$-module. Since filtered colimits are exact we thus obtain that
 		$$\underbrace{\colimit_{R'\subseteq R} A\otimes_{R'[T]} R'[T]\otimes_{R'} \free{R'_\blacksquare}{S}}_{\cong \free{(A, R)_\blacksquare}{S}} \to \underbrace{\colimit_{R'\subseteq R} A\otimes_{R'[T]} \free{R'[T]_\blacksquare}{S}}_{\cong \free{A_\blacksquare}{S}} \to \colimit_{R'\subseteq R} A\otimes_{R'[T]} Q_{R'} \to 0$$
 	is exact as well and hence that the quotient must be $Q$. Now if $s\in \laurent{R'}{T^{-1}}\subseteq A_\infty$ and $a \in A\otimes_{R''[T]} Q_{R''} \subseteq Q$ then consider equivalent $s$ in $\laurent{R'''}{T^{-1}}$ and $a$ in $A\otimes_{R'''[T]}Q_{R'''}$ with $R'''=\Z[R'\cup R'']$ still $\Z$-finitely generated. Then $s$ acts on $a$ as $A\otimes_{R'''[T]}Q_{R'''}$ is an $\laurent{R'''}{T^{-1}}$-module. This makes $Q$ into an $A_\infty = \colimit_{R'\subseteq R} \laurent{R'}{T^{-1}}$-module.
\end{proof}

%\newpage
%
%\section{The virtual case of \texorpdfstring{$\id\colon (R[T], R)_\blacksquare \to R[T]_\blacksquare$}{Identities of Solid Rings}}

\begin{lemma}
	The kernel of the completed scalar extension $j^\ast\colon \derived{(A, R)_\blacksquare} \to \derived{A_\blacksquare}$ is the full subcategory of $\derived{(A, R)_\blacksquare}$ of $A_\infty$-modules.
\end{lemma}
\begin{proof}
	Suppose that $M^\bullet$ is an $A_\infty$-module. By lemma \ref{lemma: Ainfty-modules are reflective} we have that $M^\bullet \cong A_\infty\Dotimes_{(A, R)_\blacksquare} M^\bullet$. Since $j^\ast$ is monoidal by lemma \ref{lemma: completed scalar extension is monoidal}, we obtain that
	$$j^\ast M^\bullet \cong j^\ast(A_\infty\Dotimes_{(A, R)_\blacksquare}M^\bullet)\cong j^\ast A_\infty \Dotimes_{A_\blacksquare} j^\ast M^\bullet$$
	naturally is a $j^\ast A_\infty = A_\blacksquare \Dotimes_{(A, R)_\blacksquare} A_\infty$-module. Now for any $N^\bullet \in \derived{A_\blacksquare}$ we find
	$$\eRHom_A(A_\blacksquare\Dotimes_{(A, R)_\blacksquare} A_\infty, N^\bullet)\cong \eRHom_A(A_\infty, N^\bullet) \cong 0$$
	by lemma \ref{lemma: Ainfty-modules admit not extensions by Ablack-completes} as $N^\bullet$ is $A_\blacksquare$-complete and thus $(A, R)_\blacksquare$-complete. Thus, $j^\ast M$ is the zero-ring and any module over it must be trivial. In particular, we obtain that $j^\ast M^\bullet \cong 0$ so that $M^\bullet$ is in the kernel of $j^\ast$.
	
	Now suppose that $M^\bullet \in \derived{(A, R)_\blacksquare}$ is in the kernel of $j^\ast$, \ie $j^\ast M^\bullet \cong 0$. Now $M^\bullet$ admits a left resolution $P^\bullet \to M^\bullet$ by a $K$-projective complex $P^\bullet$ that, as a complex, is a sequential colimit $P^\bullet = \colimit_{n\ge -1} P^\bullet_n$ of bounded to the right complexes $P^\bullet_n$ which are term-wise compact projectives $P^m_n = \bigoplus_{i \in I^m_n} \free{(A, R)_\blacksquare}{S_{i, n}^m}$ of $\derived{(A, R)_\blacksquare}$. Denote by $\phi\colon (A, R)_\blacksquare \to A_\blacksquare$ the map $\id_A$ of analytic rings. Then $K$-projectivity of $P^\bullet$ gives
	$$j^\ast M^\bullet \ldef \LeftD\completedScalarExtension{\phi}{A_\blacksquare}M^\bullet \isorightarrow \completedScalarExtension{\phi}{A_\blacksquare}P^\bullet = A_\blacksquare \otimes_{(A, R)_\blacksquare} P^\bullet,$$
	implying that $\completedScalarExtension{\phi}{A_\blacksquare}P^\bullet$ is acyclic as $j^\ast M^\bullet \cong 0$. For any $n\ge -1$ consider the short exact sequence
	$$0\to P_n^\bullet \to \completedScalarExtension{\phi}{A_\blacksquare}P_n^\bullet \to Q_n^\bullet \to 0$$
	where $P_n^\bullet \to \completedScalarExtension{\phi}{A_\blacksquare} P_n^\bullet$ is given by the unit of the adjunction $\completedScalarExtension{\phi}{A_\blacksquare} \ladj \scalarRestriction{\phi}$. Observe that for any $m\in \Z$ the above short exact sequence corresponds to the short exact sequence
	$$0\to \underbrace{\bigoplus_{i\in I_n^m} \free{(A, R)_\blacksquare}{S_{i, n}^m}}_{=P_n^m} \to \underbrace{\bigoplus_{i\in I_n^m} \free{A_\blacksquare}{S_{i, n}^m}}_{\cong \completedScalarExtension{\phi}{A_\blacksquare}P_n^m} \to Q_n^m \to 0,$$
	where $\completedScalarExtension{\phi}{A_\blacksquare}P_n^m \cong \bigoplus_{i\in I_n^m} \free{A_\blacksquare}{S_{i, n}^m}$ as $\completedScalarExtension{\phi}{A_\blacksquare}$ is the unique colimit preserving extension of $\free{(A, R)_\blacksquare}{S}\mapsto \free{A_\blacksquare}{S}$ by proposition \ref{proposition: scalar restriction and extension of complete modules}. Lemma \ref{lemma: special cokernels are Ainfty-modules, general case} now implies that these cokernels $Q_n^m$ are $A_\infty$-modules, yielding that the complexes $Q_n^\bullet$ and their colimit $Q^\bullet = \colimit_{n\ge -1} Q_n^\bullet$ are $A_\infty$-modules as well. Furthermore, since filtered colimits are exact, the system of short exact sequences
	$$0\to P_n^\bullet \to \completedScalarExtension{\phi}{A_\blacksquare}P_n^\bullet \to Q_n^\bullet \to 0$$
	gives rise, in the colimit, to a short exact sequence
	$$0\to P^\bullet \to \completedScalarExtension{\phi}{A_\blacksquare}P^\bullet \to Q^\bullet \to 0$$
	since $\completedScalarExtension{\phi}{A_\blacksquare}$ (as a left adjoint) commutes with colimits of complexes. This short exact sequence then gives rise to a distinguished triangle
	$$P^\bullet \to \completedScalarExtension{\phi}{A_\blacksquare}P^\bullet \to Q^\bullet \to P^\bullet[1]$$
	in $\derived{(A, R)_\blacksquare}$. Now $\completedScalarExtension{\phi}{A_\blacksquare}P^\bullet$ is acyclic, yielding (after rotating the triangle) that $Q^\bullet[-1] \cong P^\bullet$ by \cite[\href{https://stacks.math.columbia.edu/tag/05QR}{Tag 05QR}]{stacks-project}. But $Q^\bullet$ is an $A_\infty$-module, implying that $P^\bullet$ and hence $M^\bullet$ are $A_\infty$-modules as well.
\end{proof}

We can now construct for $j\colon \Spec A_\blacksquare \to \Spec (A, R)_\blacksquare$ its associated exceptional direct image functor. 
\begin{lemma}[The left adjoint of completed scalar extension]
	\label{lemma: the left adjoint of completed scalar extension}
	
	The endofunctor
		$$T\colon \derived{(A, R)_\blacksquare}\to\derived{(A, R)_\blacksquare},\;M\to M\Dotimes_{(A, R)_\blacksquare} (A_\infty/A)[-1]$$
	given by twisting with $(A_\infty/A)[-1]$ factors uniquely over a fully faithful functor
		$$j_!\colon \derived{A_\blacksquare}\to\derived{(A, R)_\blacksquare}$$
	left adjoint to $j^\ast$ such that $j_! j^\ast M = j_! A \Dotimes_{(A, R)_\blacksquare} M$ for all $M\in\derived{(A, R)_\blacksquare}$.
\end{lemma}
\begin{proof}
	Observe that if this left adjoint $j_!$ to $j^\ast$ exists, then $j^\ast\circ j_!$ will be left adjoint to $j^\ast \circ j_\ast$ which is the identity by full faithfulness of scalar restriction $j_\ast$. Thus, $j^\ast\circ j_!$ is the identity as well and $j_!$ is fully faithful.
	
	For the existence of the adjoint $j_!$ and its representation as a tensor product it suffices to see that for all $M, N \in \derived{(A, R)_\blacksquare}$ we have
		$$\eRHom_A(TM, N)\isorightarrow \eRHom_A(TM, j_\ast j^\ast N) \isoleftarrow \eRHom_A(M, j_\ast j^\ast N)$$
	in $\derived{\condensedAb}$. Then if $M$ is $A_\blacksquare$-complete and $N$ is $(A, R)_\blacksquare$-complete we obtain by lemma \ref{lemma: eRHom enriches Hom-sets in the derived category} that
		$$\Hom_{(A, R)_\blacksquare}(T j_\ast M, N)\cong \Hom_{(A, R)_\blacksquare}(j_\ast M, j_\ast j^\ast N) = \Hom_{A_\blacksquare}(M, j^\ast N)$$
	where the equality is due to $j_\ast$ (the scalar restriction along $\id_A\colon (A, R)_\blacksquare\to A_\blacksquare$) being fully faithful. Hence $j_!\ldef T\circ j_\ast \colon \derived{A_\blacksquare}\to\derived{(A, R)_\blacksquare}$ is the desired adjoint, and we find that $j_! A \cong TA \cong (A_\infty/A)[-1]$ as $j_\ast A = A$ is the unit object for $\Dotimes_{(A, R)_\blacksquare}$. In particular, $T = TA \Dotimes_{(A, R)_\blacksquare} -$ and $j_! = j_!A \Dotimes_{(A, R)_\blacksquare} j_\ast-$. Now if $M\in \derived{(A, R)_\blacksquare}$, then $j_\ast j^\ast M \cong A\Dotimes_{(A, R)_\blacksquare} M$ which implies that
		$$j_!j^\ast M \cong j_!A\Dotimes_{(A, R)_\blacksquare} j_\ast j^\ast M \cong j_!A\Dotimes_{(A, R)_\blacksquare} (A \Dotimes_{(A, R)_\blacksquare} M) \cong j_!A\Dotimes_{(A, R)_\blacksquare} M$$
	as claimed. It remains to prove the two isomorphisms of derived $\Hom$s.
	
	For the first isomorphism observe that by definition of a triangulated category, there is a distinguished triangle
	$$N \to j_\ast j^\ast N \to C \to N[1].$$
	Application of $j^\ast = A_\blacksquare\Dotimes_{(A, R)_\blacksquare} -$ then yields another distinguished triangle
		$$j^\ast N\to j^\ast j_\ast j^\ast N \to j^\ast C \to j^\ast N[1]$$
	where $j^\ast j_\ast j^\ast N \cong j^\ast N$ by full faithfulness of $j_\ast$. By \cite[\href{https://stacks.math.columbia.edu/tag/05QR}{Tag 05QR}]{stacks-project} we obtain that $j^\ast C \cong 0$ as the first morphism is an isomorphism (as it corresponds to the identity of $j^\ast N$). By lemma \ref{lemma: Ainfty-modules are reflective} this implies that $C$ is an $A_\infty$-module. By the same lemma the inclusion of $A_\infty$-modules is fully faithful (hence $\eHom$s and $\eRHom$s of the categories agree) and has a left adjoint $A_\infty\Dotimes_{(A, R)_\blacksquare} -$. This adjunction arises from deriving the scalar extension and restriction between $(A_\infty, R)_\blacksquare$-modules and $(A, R)_\blacksquare$-modules. This adjunction is $\condensedAb$-enriched adjunction and derivable by \ref{lemma: enriched adjunctions are stable under derivation} so we obtain that the original adjunction is $\derived{\condensedAb}$-enriched and hence that
		$$\eRHom_A(TM, C) \cong \eRHom_A(A_\infty\Dotimes_{(A, R)_\blacksquare} TM, C).$$
	Consider the short exact sequence $0\to A\to A_\infty\to A_\infty/A \to 0$ and apply $A_\infty\Dotimes_{(A, R)_\blacksquare} -$ to obtain the distinguished triangle
		$$A_\infty\Dotimes_{(A, R)_\blacksquare} A \to A_\infty\Dotimes_{(A, R)_\blacksquare}A_\infty \to A_\infty\Dotimes_{(A, R)_\blacksquare} (A_\infty/A) \to (\cdots)[1].$$
	But $A$ is the unit of $\Dotimes_{(A, R)_\blacksquare}$ and $A_\infty$ is idempotent, implying that also
	$$A_\infty\xrightarrow{\id} A_\infty \to A_\infty\Dotimes_{(A, R)_\blacksquare} (A_\infty/A) \to (\cdots)[1]$$
	is distinguished. By \cite[\href{https://stacks.math.columbia.edu/tag/05QR}{Tag 05QR}]{stacks-project} we thus have $A_\infty\Dotimes_{(A, R)_\blacksquare} (A_\infty/A)\cong 0$ as the first morphism is an isomorphism. This implies that also $A_\infty\Dotimes_{(A, R)_\blacksquare} TM = A_\infty\Dotimes_{(A, R)_\blacksquare} (A_\infty/A) \Dotimes_{(A, R)_\blacksquare} M[-1] \cong 0$ and hence also $\eRHom_A(A_\infty\Dotimes_{(A, R)_\blacksquare} T M, C) \cong 0$ giving $\eRHom_A(TM, C) \cong 0$ as well. Thus application of $\eRHom_A(TM, -)$ to the initial triangle then yields that also
		$$\eRHom_A(TM, N)\to\eRHom_A(TM, j_\ast j^\ast N) \to 0 \to (\cdots)[1]$$
	is distinguished. Now by \cite[\href{https://stacks.math.columbia.edu/tag/05QR}{Tag 05QR}]{stacks-project} the remaining non-trivial morphism
		$$\eRHom_A(TM, N)\to\eRHom_A(TM, j_\ast j^\ast N)$$
	is an isomorphism, as claimed.
	
	
	For the second isomorphism consider again $0\to A\to A_\infty\to A_\infty/A \to 0$ which induces a distinguished triangle
	$$A\Dotimes_{(A, R)_\blacksquare} M\to A_\infty\Dotimes_{(A, R)_\blacksquare} M \to (A_\infty/A) \Dotimes_{(A, R)_\blacksquare} M \to (\cdots)[1].$$
	By observing that $A$ is the unit for $\Dotimes_{(A, R)_\blacksquare}$ and by rotating the triangle we obtain that
	$$(A_\infty/A) \Dotimes_{(A, R)_\blacksquare} M[-1] \to M \to A_\infty\Dotimes_{(A, R)_\blacksquare} M \to (A_\infty/A) \Dotimes_{(A, R)_\blacksquare} M$$
	is distinguished as well. This is precisely
	\begin{equation}
		\label{equation: distinguished triangle for T M}
		T M \to M \to A_\infty\Dotimes_{(A, R)_\blacksquare} M \to T M[1].
	\end{equation}
	Clearly $A_\infty\Dotimes_{(A, R)_\blacksquare} M$ is an $A_\infty$-module implying that
	$$\eRHom_A(A_\infty\Dotimes_{(A, R)_\blacksquare} M, j_\ast j^\ast N)\cong 0$$
	by lemma \ref{lemma: Ainfty-modules admit not extensions by Ablack-completes}, as the $A$-module $j_\ast j^\ast N = A_\blacksquare\Dotimes_{(A, R)_\blacksquare} N$ is $A_\blacksquare$-complete. Thus, application of $\eRHom_A(-, j_\ast j^\ast N)$ 
	to \ref{equation: distinguished triangle for T M} yields another distinguished triangle in which one term vanishes, implying that the remaining non-trivial morphism
		$$\eRHom_A(M, j_\ast j^\ast  N)\to \eRHom_A(T M, j_\ast j^\ast N)$$
	is an isomorphism, as seen by rotating the triangle and applying \cite[\href{https://stacks.math.columbia.edu/tag/05QR}{Tag 05QR}]{stacks-project}.
\end{proof}


%\begin{lemma}
%	\begin{todo}
	%		 Generalization
	%	\end{todo}
%	Suppose that the $\Z$-algebra $A$ is finitely generated. For any set $I$ the natural map
%	$$j_! \prod_I A \to \prod_I (A_\infty/A)[-1]$$
%	is an isomorphism of $A$-modules and the right-hand side is a compact object of $\derived{R_\blacksquare}$.
%\end{lemma}
%\begin{proof}
%	\begin{todo}
	%		proof
	%	\end{todo}
%\end{proof}


\begin{lemma}[$j_!$ and compact generators]
	\label{lemma: j lower shriek and compact generators}
	
	Denote by $\phi\colon R\to R[T]$ the map of discrete rings. For any extremally disconnected $S$ the scalar restriction $(\scalarRestriction{\phi} \circ j_!)\free{A_\blacksquare}{S}$ is a derived compact object of $\derived{R_\blacksquare}$.
\end{lemma}
\begin{proof}
	Observe that
		$$j_!\free{A_\blacksquare}{S}\cong j_!j^\ast\free{(A, R)_\blacksquare}{S}\cong j_!A\Dotimes_{(A, R)_\blacksquare} \free{(A, R)_\blacksquare}{S}$$
	where the last isomorphism follows from lemma \ref{lemma: the left adjoint of completed scalar extension}. Now in \cite[Observation 8.12, p.57]{condenseddotpdf} Clausen and Scholze reduce as follows:
		$$\scalarRestriction{\phi}(j_!A \Dotimes_{(A, R)_\blacksquare} \free{(A, R)_\blacksquare}{S}) \cong (\scalarRestriction{\phi}j_!A) \Dotimes_{R_\blacksquare} \free{R_\blacksquare}{S}.$$
	The author was not able to reproduce this result, neither in this general case nor in case that $R$ is $\Z$-finitely generated (they claim this isomorphism for $R=\Z$). We thus have to make this an \assumptionPreWord. Now $j_!A \cong (A_\infty/A)[-1]$ and
		$$A_\infty/A \cong \colimit_{R'\subseteq R} T^{-1}\cdot \series{R'}{T^{-1}} \subseteq T^{-1}\cdot \series{R}{T}$$
	as for each $R'$ we have $\laurent{R'}{T^{-1}}/R'[T] \cong T^{-1}\cdot \series{R'}{T^{-1}}$. This then yields that
		$$A_\infty/A \cong \colimit_{R'\subseteq R} R\otimes_{R'} T^{-1}\cdot\series{R'}{T^{-1}} \cong \colimit_{R'\subseteq R} R\otimes_{R'}\prod_\N R'$$
	so that $A_\infty/A$ is a compact object of $\mod{R_\blacksquare}$ by lemma \ref{lemma: more compact projectives}. Thus, by corollary \ref{corollary: tensor products of compacts, general case} the tensor product $(A_\infty/A)\Dotimes_{R_\blacksquare}\free{R_\blacksquare}{S}$ is a compact object of $\mod{R_\blacksquare}$ and we finally find that the complex $(\scalarRestriction{\phi}j_! \free{A_\blacksquare}{S}) = (A_\infty/A)\Dotimes_{R_\blacksquare}\free{R_\blacksquare}{S}[-1]$ is a derived compact of $\derived{R_\blacksquare}$.
	
	
%Since $\derived{A_\blacksquare}$ is generated by the derived compact objects $\free{A_\blacksquare}{S}$ for $S$ extremally disconnected and since $j_!$ commutes with direct sums it is enough to prove the claim on these generators. Thus let $S$ be extremally disconnected and
%	Now $j_! A \cong A_\infty/A[-1]$ and $A[-1] \to A_\infty[-1]\to A_\infty/A[-1] \to A$ being distinguished imply that $j_!\free{A_\blacksquare}{S} \cong j_!A\Dotimes_{(A, R)_\blacksquare} \free{(A, R)_\blacksquare}{S}$ is derived compact if both $A[-1]\Dotimes_{(A, R)_\blacksquare} \free{(A, R)_\blacksquare}{S}$ and $A_\infty[-1]\Dotimes_{(A, R)_\blacksquare} \free{(A, R)_\blacksquare}{S}$ are derived compact (it is straightforward to see that if $2$ out of $3$ legs of a distinguished triangle are derived compact then so is the third). As $A$ is the unit object of $\Dotimes_{(A, R)_\blacksquare}$ it follows that the first tensor product is simply $\free{(A, R)_\blacksquare}{S}$ which certainly is derived compact. It remains to see that (the shift of) the second tensor product $A_\infty\Dotimes_{(A, R)_\blacksquare} \free{(A, R)_\blacksquare}{S}$ is derived compact. But
%	$$A_\infty \cong \colimit_{R'\subseteq R} A\otimes_{R'}\laurent{R'}{U} \cong A\otimes_R \colimit_{R'\subseteq R} R\otimes_{R'} \prod_\N R'$$
%	and similarly
%	$$\free{(A, R)_\blacksquare}{S} \cong A\otimes_R \colimit_{R'\subseteq R} R\otimes_{R'} \prod_I R'$$
%	where $I$ is determined by Specker's theorem \ref{theorem: specker's theorem} yielding that 
%	$$A_\infty\Dotimes_{(A, R)_\blacksquare} \free{(A, R)_\blacksquare}{S} \cong A\otimes_R \colimit_{R'\subseteq R} R\otimes_{R'} \prod_{\N\times I} R'$$
%	analogous to the proof of corollary \ref{corollary: tensor products of compacts, general case} and the further that this $A$-module is derived compact by (the proof of) lemma \ref{lemma: more compact projectives} (as the retraction is preserved by the extension $A\otimes_{R} -$).
	
	
	
	%	 By lemma \ref{lemma: presentation for Afinty, general case} there is a short exact sequence
	%		$$0\to\colimit_{R'\subseteq R} \series{R'}{U}[T] \to \colimit_{R'\subseteq R} \series{R'}{U}[T] \to A_\infty \to 0$$
	%	with the first two terms compact projective. Writing
	%		$$\colimit_{R'\subseteq R} \series{R'}{U}[T] \cong \colimit_{R'\subseteq R} A\otimes_R R\otimes_{R'}\series{R'}{U}\cong  A\otimes_R \mleft(\colimit_{R'\subseteq R} R\otimes_{R'}\series{R'}{U}[T]\mright)\rdef A\otimes_R Z$$
	%	we thus find that $A_\infty$ admits the projective left resolution $0\to A\otimes_R Z\to A \otimes_R Z\to 0$ which is the derived completed scalar extension $\LeftD\completedScalarExtension{\phi}{(A, R)_\blacksquare}(0\to Z\to \Z \to 0)$.
	%	Since derived completed scalar extension $\completedScalarExtension{\phi}{(A, R)_\blacksquare}$ is strong symmetric monoidal by lemma \ref{lemma: completed scalar extension is monoidal} and preserves derived compactness we thus might as well prove that
	%		$$(0\to Z\to Z\to 0)\Dotimes_{R_\blacksquare} \free{R_\blacksquare}{S}$$
	%	is derived compact. As
	%		$$(0\to Z\to Z\to 0)\Dotimes_{R_\blacksquare} \free{R_\blacksquare}{S} \cong (0\to Z\otimes_{R_\blacksquare} \free{R_\blacksquare}{S} \to Z\otimes_{R_\blacksquare} \free{R_\blacksquare}{S} \to 0)$$
	%	this further reduces to the claim that $Z\otimes_{R_\blacksquare} \free{R_\blacksquare}{S}$ is a compact projective object of $\mod{R_\blacksquare}$. Using a cofinality argument (to simplify the two colimits appearing in the definition of $Z$ and $\free{R_\blacksquare}{S}$ to a single colimit), that $\completion{-}{R_\blacksquare}$ is strong monoidal and that tensor products commute with arbitrary colimits, a long but straightforward calculation shows that
	%		$$Z\otimes_{R_\blacksquare}\free{R_\blacksquare}{S}\cong \colimit_{R'\subseteq R} R\otimes_{R'} (\series{R'}{U}\otimes_{R'_\blacksquare} \free{R'_\blacksquare}{S}).$$
	%	By Specker's theorem \ref{theorem: specker's theorem} there is a set $I$ such that for each such $R'$ we have $\free{R'_\blacksquare}{S}\cong \prod_I R'$. Similarly we find that $\series{R'}{U}\cong \prod_\N R'$.  Akin to proposition \ref{proposition: tensor products of compacts} we thus find that $\series{R'}{U}\otimes_{R'_\blacksquare} \free{R'_\blacksquare}{S} \cong \prod_{\N\times I} R'$. Choosing $S'$ extremally disconnected large enough such that the $I'$ associated to $S'$ by Specker satisfies $\vert I'\vert \ge \vert I\times \N\vert$ we find that $\series{R'}{U}\otimes_{R'_\blacksquare}\free{R'_\blacksquare}{S}\cong \prod_{I\times \N} R'$ is a retract of $\free{R'_\blacksquare}{S'}\cong\prod_{I'} R'$. We obtain that
	%		$$Z\otimes_{R_\blacksquare} \free{R_\blacksquare}{S} \cong \colimit_{R'\subseteq R} R\otimes_{R'}\series{R'}{U}\otimes_{R'_\blacksquare}\free{R'_\blacksquare}{S}$$
	%	is a retract of
	%		$$\colimit_{R'\subseteq R} R\otimes_{R'}\free{R'_\blacksquare}{S'}\rdef \free{R_\blacksquare}{S'}$$
	%	which is a compact projective! As the retract of any compact projective is compact projective again we obtain the claim.
\end{proof}

%\begin{corollary}
%	For every compact generator $\free{A_\blacksquare}{S}$ of $\derived{A_\blacksquare}$ with $S$ extremally disconnected the object $j_! \free{A_\blacksquare}{S} \in \derived{(A, R)_\blacksquare}$ is compact in $R_\blacksquare$.
%\end{corollary}
%
%\begin{proof}
%	Indeed, for any $S$ extremally disconnected we have
%	$$j_!\free{A_\blacksquare}{S} \cong A_\infty/A\Dotimes_{(A, R)_\blacksquare} \free{(A, R)_\blacksquare}{S}[-1]$$
%	by the proof of said lemma. But $(A, R)_\blacksquare$-completion is (strong) monoidal by lemma \ref{lemma: completed scalar extension is monoidal} showing that $j_!\free{A_\blacksquare}{S}$ (up to shift) is the completion of
%	$$A_\infty/A\Dotimes_A \free{(A, R)_\blacksquare}{S}\cong A_\infty/A \Dotimes_R \free{R_\blacksquare}{S}.$$
%	Now by Specker's theorem there is some set $I$ such that $\free{R_\blacksquare}{S} \cong \colimit_{R'\subseteq R} R\otimes_{R'}\prod_I R'$. We obtain a similar representation of $A_\infty/A$. Indeed, $A_\infty\cong \colimit_{R'\subseteq R} \laurent{R'}{T^{-1}}$ and $A = R[T] \cong \colimit_{R'\subseteq R} R'[T]$ give that
%	$$A_\infty/A\cong\colimit_{R'\subseteq R} \laurent{R'}{T^{-1}}/R'[T] \cong \colimit_{R'\subseteq R} T^{-1}\cdot\series{R'}{T^{-1}}.$$
%	Thus $A_\infty/A \cong \colimit_{R'\subseteq R} R\otimes_{R'}\prod_\N R'$ and hence corollary \ref{corollary: tensor products of compacts, general case} we find that
%	$$A_\infty/A\Dotimes_{R_\blacksquare} \free{R_\blacksquare}{S} \cong \colimit_{R'\subseteq R} R\otimes_{R'}\prod_{\N\times I} R'$$
%	is compact.
%	
%	
%	That $\scalarRestriction{\phi}$ preserves derived compactness follows analogous to the proof of lemma \ref{lemma: scalar restriction preserves compactness for finite Tor-dimension} after realizing that we do not need finiteness of $R\to R[X]$ as $P^\bullet = R[X][0]$ is already a(n infinite) free resolution and that also infinite direct sums of compact objects are compact???
%\end{proof}


\section{Exceptional direct images of virtual maps}

Using the most basic virtual map $\Spec R[T]_\blacksquare \to \Spec (R[T], R)_\blacksquare$ associated to the affine space $\A_R^1$ one can construct the exceptional direct image for many other virtual maps too.
\begin{proposition}
	\label{proposition: lower shriek for virtual maps}
	
	Let $R\to S\to A$ be morphisms of discrete rings. The scalar restriction of the morphism of analytic rings $\id_A\colon (A, R)_\blacksquare\to (A, S)_\blacksquare$ is a functor
	$$j_\ast\colon \derived{(A, S)_\blacksquare} \to \derived{(A, R)_\blacksquare}$$
	that admits a left adjoint $j^\ast \ldef (A, S)_\blacksquare \otimes^\LeftD_{(A, R)_\blacksquare}  -$ which in turn admits a left adjoint
	$$j_!\colon\derived{(A, S)_\blacksquare} \to \derived{(A, R)_\blacksquare}.$$
	Then $j_!$ and $j^\ast$ satisfy
	$$j_!j^\ast M = j_! A  \Dotimes_{(A, R)_\blacksquare} M$$
	for all $M\in \derived{(A, R)_\blacksquare}$.
\end{proposition}
\begin{proof}
	Since $(A, R)_\blacksquare$ and $(A, S)_\blacksquare$ are analytic by corollary \ref{corollary: analyticity of (A, R)black}, the map of analytic rings $\id_A\colon (A, R)_\blacksquare\to(A, S)_\blacksquare$ induces by proposition \ref{proposition: scalar restriction and extension of complete modules} the desired functors $j_\ast$ and $j^\ast$, the scalar restriction and (completed) scalar extension of derived modules respectively.
	
	Notice that if the statement is true for $R\to S \to S$ then statement is true for $R \to S \to A$ in general. To see this, denote by $j\colon \Spec S_\blacksquare \to \Spec (S, R)_\blacksquare$ and $k\colon \Spec (A, S)_\blacksquare \to \Spec (A, R)_\blacksquare$ the virtual maps and assume $j_!$ exists and satisfies the desired properties. Consider the following diagram
	$$\begin{tikzcd}
		\derived{S_\blacksquare} &\derived{(S, R)_\blacksquare}\\
		\derived{(A, S)_\blacksquare} &\derived{(A, R)_\blacksquare}
		\ar[from=1-1, to=1-2, shift left=0.2em, "j_!"]
		\ar[from=1-2, to=1-1, shift left=0.2em, "j^\ast"]
		\ar[from=2-1, to=2-2, shift left=0.2em, "k_!", dotted]
		\ar[from=2-2, to=2-1, shift left=0.2em, "k^\ast"]
		\ar[from=1-1, to=2-1, shift right=0.2em, "e", swap]
		\ar[from=2-1, to=1-1, shift right=0.2em, "r", swap]
		\ar[from=1-2, to=2-2, shift right=0.2em, "e'", swap]
		\ar[from=2-2, to=1-2, shift right=0.2em, "r'", swap]
	\end{tikzcd}$$
	where $e\ladj r$ and $e'\ladj r'$ are the respective (derived completed) scalar extension and restriction along $S\to A$. Now define $k_!$ to be the unique colimit preserving extension of $\free{(A, S)_\blacksquare}{T} \mapsto (e'\circ j_!)\free{S_\blacksquare}{T}$ where $T$ is extremally disconnected. Then for all such $T$ and any $(A, S)_\blacksquare$-complete $M$ we have
	\begin{align*}
		\Hom_A(k_!\free{(A, S)_\blacksquare}{T}, M) &\ldef \Hom_A((e'\circ j_!)\free{S_\blacksquare}{T}, M)\\
		&\cong \Hom_S(j_!\free{S_\blacksquare}{T}, r'(M))\\
		&\cong \Hom_S(\free{S_\blacksquare}{T}, (j^\ast \circ r')(M))\\
		&\cong \Hom_S(\free{S_\blacksquare}{T}, (r\circ k^\ast)(M))\\
		&\cong \Hom_A(e\free{S_\blacksquare}{T}, k^\ast (M)) \rdef \Hom_A(\free{(A, R)_\blacksquare}{T}, k^\ast(M))
	\end{align*}
	proving the adjunction on generators, implying that $k_! \ladj k^\ast$ in general. Now a straightforward but long calculation yields that on generators
		$$k_!k^\ast \free{(A, R)_\blacksquare}{T} \cong k_! A \Dotimes_{(A, R)_\blacksquare} \free{(A, R)_\blacksquare}{T}$$
	using that $k_!A \cong k_!\free{(A, R)_\blacksquare}{\ast}$ and that the (derived completed) scalar extension $e'$ is (strong) symmetric monoidal by lemma \ref{lemma: completed scalar extension is monoidal}. As both $k_!$ and $k^\ast$ commute with arbitrary colimits (as left adjoints) this gives the desired representation. This proves the claimed reduction.

	We can thus assume without loss of generality that $S\to A$ is surjective as this certainly covers the case of $S=A$. Since any surjection $S\to A$ is finite by \cite[\href{https://stacks.math.columbia.edu/tag/04VT}{Tag 04VT}]{stacks-project} we obtain by lemma \ref{lemma: analytic rings associated to finite morphisms} that $\id_A\colon (A, S)_\blacksquare \to A_\blacksquare$ is an isomorphism. Hence, $(A, S)_\blacksquare$ is independent of this surjection and we can assume further that $S = R[T_1, ..., T_n]$. Reducing again $R\to S\to A$ to $R\to S\to S$ we can assume both rings $S = A = R[T_1, \dots, T_n]$ to be polynomial algebras over $R$. If $n=0$ there is nothing to show. If $n\ge 2$ then the map $\id_A \colon (A, R)_\blacksquare \to (A, S)_\blacksquare$ factors as
	$$(A, R)_\blacksquare \to (A, R[T_1, \dots, T_{n-1}])_\blacksquare \to (A, S)_\blacksquare$$
	which inductively reduces the statement to the case of $n=1$. Indeed if we denote the associated virtual morphisms by $k\colon \Spec(A, R[T_1, \dots, T_{n-1}])_\blacksquare\to \Spec (A, R)_\blacksquare$ and $\ell\colon \Spec (A, S)_\blacksquare \to \Spec (A, R[T_1,\dots, T_{n-1}])_\blacksquare$ and assume $k_!$ and $\ell_!$ to satisfy the claimed properties then surely $(k\circ \ell)_!\ldef k_!\circ \ell_!$ is left adjoint to $j^\ast = (k\circ \ell)^\ast \cong \ell^\ast \circ k^\ast$ and a straightforward calculation yields that $(k\circ \ell)_!\circ (k\circ \ell)^\ast \cong (k\circ\ell)_! A \otimes_{(A, R)_\blacksquare} -$ after observing that $k_!$ is (strong) symmetric monoidal by \ref{lemma: completed scalar extension is monoidal}. Hence, the claim finally reduces to the case of $A = S = R[T]$ being a univariate polynomial algebra. But this is precisely lemma \ref{lemma: the left adjoint of completed scalar extension}.
\end{proof}


\section{The Duality}

Recall that there are several notions of \emph{dimension} of a module over a ring. Two of them will prove relevant to us.
\begin{definition}[Tor-dimension \& projective dimension] 
	Let $R$ be a discrete ring and $M$ a discrete $R$-module. Call the smallest length of a (possibly infinite) left resolution of $M$ by discrete flat $R$-modules the \emph{Tor-dimension} of $M$. Similarly, call the smallest length of a (possibly infinite) left resolution of $M$ by discrete projective $R$-modules the \emph{projective dimension} of $M$.
	
	If $f\colon \Spec A\to \Spec R$ is a morphism of affine schemes then one calls $f$ of finite Tor-dimension respectively of finite projective dimension if the $A$-module $R$ has finite Tor-dimension respectively projective dimension.
\end{definition}

As any projective resolution of $M$ is also flat, we find that the projective dimension is never greater than the Tor dimension. When do they agree?
\begin{lemma}[Agreement of dimensions, {\cite[Prop. 4.1.5]{weibel-homological-algebra}}]
	\label{lemma: agreement of dimensions}
	
	If the discrete ring $R$ is Noetherian and $M$ is a finitely generated discrete $R$-module then the Tor-dimension and the projective dimension of $M$ agree.
\end{lemma}

We obtain the following useful result.
\begin{lemma}
	\label{lemma: scalar restriction preserves compactness for finite Tor-dimension}
	
	Let $f\colon \Spec A\to\Spec R$ be a finite morphism of finite Tor-dimension and let $R$ be Noetherian. Denote by $\phi\colon R\to A$ the corresponding morphism of rings. Then scalar restriction $\scalarRestriction{\phi}\colon \derived{(A, R)_\blacksquare}\to \derived{R_\blacksquare}$ preserves derived compactness.
\end{lemma}
\begin{proof}
	As $\derived{(A, R)_\blacksquare}$ is generated by the compact objects $\free{(A, R)_\blacksquare}{S}$ for $S$ extremally disconnected we only need to show that $\scalarRestriction{\phi}$ preserves compactness on those generators. So let $S$ be extremally disconnected.
	
	Since $f$ is finite we know that $A$ is a finitely generated $R$-module. As $f$ is of finite Tor-dimension we thus know by lemma \ref{lemma: agreement of dimensions} that the $R$-module $A$ has finite projective dimension. Choose a bounded left resolution $P^\bullet$ of $A$ by discrete finitely generated projective $R$-modules. Write $P^\bullet$ as a direct summand of a bounded complex $F^\bullet$ of discrete finitely generated free $R$-modules. Then $P^\bullet \otimes_{R_\blacksquare} \free{R_\blacksquare}{S}$ is a direct summand of the complex $F^\bullet \otimes_{R_\blacksquare} \free{R_\blacksquare}{S}$. But in degree $n \in \Z$ there is some $m\in \N$ such that $F^n\cong R^m$ yielding that $F^n \otimes_{R_\blacksquare} \free{R_\blacksquare}{S} \cong \oplus_m \free{R_\blacksquare}{S} \cong \free{R_\blacksquare}{\sqcup_m S}$ and hence that $F^\bullet \otimes_{R_\blacksquare}{S}$ is term-wise compact and hence derived compact. Thus $P^\bullet \otimes_{R_\blacksquare} \free{A_\blacksquare}{S}$ is a direct summand of a derived compact complex hence derived compact itself.
\end{proof}

We now have everything necessary to state and prove the main theorem of this thesis.
\begin{theorem}[Exceptional direct and inverse image functors, the affine case]
	\label{theorem: exceptional direct and inverse image functors, the affine case}
	
	Let $f\colon \Spec A\to\Spec R$ be a finite type morphism of locally Noetherian affine schemes. Denote by $\phi\colon R\to A$ the corresponding morphism of rings.
	\begin{itemize}
		\item
		The map of analytic rings $R_\blacksquare\to (A, R)_\blacksquare \to A_\blacksquare$ with underlying map of rings $R\xrightarrow{\phi} A\xrightarrow{\id}A$ induces a functor
		$$f_!\colon \derived{A_\blacksquare} \xrightarrow{j_!} \derived{(A, R)_\blacksquare}\xrightarrow{\scalarRestriction{\phi}} \derived{R_\blacksquare}$$
		which commutes with all direct sums and satisfies the \emph{projection formula}
		$$f_!(N \otimes^\LeftD_{A_\blacksquare} (A_\blacksquare\otimes^\LeftD_{R_\blacksquare} M)) \cong f_!N\otimes^\LeftD_{R_\blacksquare} M$$
		for all $M\in D(R_\blacksquare)$ and all $N\in D(A_\blacksquare)$. If $f$ is of finite Tor-dimension, then $f_!$ preserves compact objects. Furthermore, $(g\circ f)_! \cong g_!\circ f_!$ if $g$ is another finite type morphism of locally Noetherian affine schemes.
		
		\item
		The functor $f_!$ admits a right adjoint
		$$f^!\colon D(R_\blacksquare)\to D(A_\blacksquare).$$
		The object $f^!R$ is a discrete and bounded to the left complex of finitely generated $A$-modules. If $f$ is of finite Tor dimension, then $f^!R$ is bounded, $f^!$ commutes with all direct sums and is given by
			$$f^!M =  f^!R \Dotimes_{A_\blacksquare} (A_\blacksquare \Dotimes_{R_\blacksquare} M)$$
		for $M\in\derived{R_\blacksquare}$, \ie $f^!$ is a \emph{twist} of scalar extension. If $f$ is a complete intersection, then $f^!R\in D(A)$ is an invertible object, \ie locally a line bundle concentrated in some degree. Furthermore, $(g\circ f)^! \iso f^!\circ g^!$ if $g$ is any other finite type morphism of locally Noetherian affine schemes.
	\end{itemize}
\end{theorem}
\begin{proof}
	Suppose that $g\colon \Spec R\to\Spec R'$ is any other morphism of affine schemes corresponding to the morphism of rings $\psi\colon R'\to R$. Observe that we have the following diagram
	$$\begin{tikzcd}
		\derived{A_\blacksquare}\ar[d, "{j^f_!}"]\ar[dr, "{j_!^{g\circ f}}"]&&\\
		\derived{(A, R)_\blacksquare}\ar[d, "\scalarRestriction{\phi}"]\ar[r] &\derived{(A, R')_\blacksquare}\ar[d, "\scalarRestriction{\phi}", swap]\ar[dr, "\scalarRestriction{(\phi\circ\psi)}"]&\\
		\derived{R_\blacksquare}\ar[r, "j^g_!"] &\derived{(R, R')_\blacksquare}\ar[r, "\scalarRestriction{\psi}"] &\derived{R'_\blacksquare}
	\end{tikzcd}$$
	where the left edge is $f_!$, the bottom edge is $g_!$ and the diagonal is $(g\circ f)_!$. If we can show that the embedded square commutes (up to natural isomorphism), then the whole diagram commutes (up to natural isomorphism) and we obtain that $(g\circ f)_!\cong g_!\circ f_!$. Similar to the proof of \ref{proposition: lower shriek for virtual maps} this follows since, in the words of Clausen and Scholze, \emphquote{... this is formal as the functor on top is simply carrying around an additional A-module structure.}
	
	By the extended adjoint functor theorem from \assumptionPreWord \ref{assumption: another adjoint functor theorem} the existence of $f^!$ as the right adjoint of $f_!$ follows from $f_!$ commuting with direct sums. Indeed, $f_! = \scalarRestriction{\phi} \circ j_!$ is the composition of two left adjoints and hence commutes with arbitrary direct sums. Using this characterization it is also clear that $(-)^!$ is compatible with composition as we already have shown $(-)_!$ to be. Furthermore, once we have shown that $f_!$ preserves compactness if $f$ is of finite Tor-dimension, then we already know from proposition \ref{proposition: preserves compact iff right adjoint preserves coproducts} that $f^!$ must commute with direct sums.
	
	Now that $(-)_!$ and $(-)^!$ are compatible with composition we can, analogous to the proof of proposition \ref{proposition: lower shriek for virtual maps}, reduce the remaining claims about $f_!$ and $f^!$ to the case that either $A$ is a free $R$-algebra or that $R\to A$ is surjective (which in particular covers the case of $f$ being a complete intersection). Indeed, since $A$ is of finite type over $R$ we find that $\phi\colon R\to A$ can be factorized as
		$$R \to R[X_1, \dots, X_n] \twoheadrightarrow A$$
	for a suitable $n$. Denote the resulting morphisms of schemes by $i\colon \Spec A\to \A_R^n$ and $\pi\colon \A_R^n \to \Spec R$ and abbreviate $R[\mathbf{X}] \ldef R[X_1,\dots, X_n]$. Clearly $R\to R[\mathbf{X}]$ is of finite Tor-dimension as the polynomial ring is $R$-free. By \cite[\href{https://stacks.math.columbia.edu/tag/0B66}{Tag 0B66}]{stacks-project} if $R\to A$ is of finite Tor-dimension then so is $R[\mathbf{X}]\twoheadrightarrow A$. Now assume all claims are true for $i$ and $\pi$, we will show that they are true for $f = \pi\circ i$ as well. Using that $A_\blacksquare\Dotimes_{R_\blacksquare} - \cong A_\blacksquare \Dotimes_{R[\mathbf{X}]_\blacksquare}(R[\mathbf{X}]_\blacksquare\Dotimes_{R_\blacksquare} -) $ one easily verifies that $f_!\cong\pi_!\circ i_!$ satisfies the projection formula. Now since $\pi$ is of finite Tor-dimension we have that $\pi^! R$ is a bounded complex of finitely generated discrete $R[\mathbf{X}]$-modules. By assumption \ref{assumption: derived adjunctions are enriched} we obtain that the adjunction $i_!\ladj i^!$ is $\condensedAb$-enriched so that
		$$\eRHom_{R[\mathbf{X}]}(i_!A, \pi^!R)\cong \eRHom_A(A, i^!\pi^!R).$$
	This yields that
		$$i^!\pi^!R \cong \eRHom_{R[\mathbf{X}]}(i_!A, \pi^!R)$$
	as $A$-modules since $\eRHom_{R[\mathbf{X}]}(i_!A, \pi^!R)$ naturally carries an $A$-module structure and since $\eRHom_A$ is the underlying (complex of) condensed abelian group(s) of $\intRHom_A$ giving that $\eRHom_A(A, i^!\pi^! R) \cong i^!\pi^!R$ as condensed abelian groups and hence as $A$-modules. As $i$ corresponds to the surjection of rings $R[\mathbf{X}]\to A$ we know by lemma \ref{lemma: analytic rings associated to finite morphisms} that the associated $j_!\colon A_\blacksquare\to (A, R[\mathbf{X}])_\blacksquare$ is simply the identity. Hence, $i_!$ is nothing more that scalar restriction and hence
		$$i^!\pi^!R\cong \eRHom_{R[\mathbf{X}]}(A, \pi^!R).$$
	Now $A$ is $R[\mathbf{X}]$-finitely generated so by Noetherianness of $R[\mathbf{X}]$ we can choose a left resolution of $A$ by discrete finitely generated projective $R[\mathbf{X}]$-modules. Replacing $A$ by its resolution we find that $\eRHom_{R[\mathbf{X}]}\cong \eHom_{R[\mathbf{X}]}^\bullet$ and thus that $i^!\pi^!R$ is a left bounded complex of discrete finitely generated $A$-modules. Now assume that $f$ is of finite Tor-dimension. Then so are $i$ and $\pi$ and hence $f_! = \pi_!\circ i_!$ preserves compactness as $f_!$ is the composition of two functors that do so too. Furthermore, as $i$ is of finite Tor-dimension we know that $i^!$ is a twist and hence we obtain that
		$$f^!R \cong i^!\pi^! R \cong i^!R[\mathbf{X}] \Dotimes_{A_\blacksquare} (A_\blacksquare \Dotimes_{R[\mathbf{X}]_\blacksquare} \pi^!R)$$
	where both $i^!R[\mathbf{X}]$ and $\pi^!R$ are bounded complexes of discrete finitely generated modules. In particular, the ordinary scalar extension of $\pi^! R$ is already $A_\blacksquare$-complete and so is the underlying $A$-balanced tensor product. Hence, in this case
		$$f^!R\cong i^!R[\mathbf{X}]\Dotimes_A(A\Dotimes_R \pi^!R)\cong i^!R[\mathbf{X}]\Dotimes_R \pi^!R$$
	is certainly a bounded complex of discrete finitely generated $A$-modules as well. It remains to see that $f^!$ is a twist.	Using the above presentation of $f^!R$, the isomorphism $A_\blacksquare\Dotimes_{R_\blacksquare} - \cong A_\blacksquare \Dotimes_{R[\mathbf{X}]_\blacksquare}(R[\mathbf{X}]_\blacksquare\Dotimes_{R_\blacksquare} -) $ of scalar extensions, that $A_\blacksquare\Dotimes_{R[\mathbf{X}]_\blacksquare} -$ is (strong) symmetric monoidal by lemma \ref{lemma: completed scalar extension is monoidal} and the twist representation of both $i^!$ and $\pi^!$ we find after a lengthy calculation that indeed $f^! M \cong (i^!\pi^!R)\Dotimes_{A_\blacksquare} (A_\blacksquare\Dotimes_{R_\blacksquare} M)$ for any $M\in\derived{R_\blacksquare}$. 
	
	This shows that we indeed can reduce.
	
	\textbf{The case of $A=R[X_1, \dots, X_n]$:}\\
	The case where $A=R[X_1,\dots, X_n]$ is a free $R$-algebra in $n\ge 2$ variables can further be reduced to the case where $n=1$ by replacing $R$ with $R=[X_1,\dots,X_{n-1}]$. Thus assume that $A=R[T]$ is a polynomial algebra in one variable.
	
	We will first verify the projection formula. Let $M\in \derived{R_\blacksquare}$ and $N\in\derived{A_\blacksquare}$. For proving the projection formula in this case of $A=R[T]$, Clausen and Scholze claim on \cite[p.57]{condenseddotpdf} that \quote{it is enough to prove the finer assertion}
	$$j_!(N\Dotimes_{A_\blacksquare} (A_\blacksquare\Dotimes_{R_\blacksquare} M)) \cong j_!N\Dotimes_{(A, R)_\blacksquare} j_\ast(A_\blacksquare \Dotimes_{R_\blacksquare} M).$$
	The author was not able to show the projection formula from the finer assertion. We thus have to make this an \assumptionPreWord. Denote by $C$ the cone of the natural morphism. We will show it to be acyclic, proving that the morphism is actually an isomorphism. Using that $j^\ast$ is (strong) symmetric monoidal by lemma \ref{lemma: completed scalar extension is monoidal} and that $j^\ast j_! \isorightarrow \id$ and $j^\ast j_\ast\isorightarrow \id$ by full faithfulness of $j_!$ and $j_\ast$ a short calculation yields that $j^\ast C \cong 0$. Furthermore, by lemma \ref{lemma: the left adjoint of completed scalar extension} $j_!$ is twisting by $j_! A$. But as shown in the proof of the same lemma we have $A_\infty\Dotimes_{(A, R)_\blacksquare} j_! A \cong 0$. Tensoring the defining distinguished triangle of $C$ we find that the images of the two objects in question are tensor products involving $A_\infty\Dotimes_{(A, R)_\blacksquare} j_! A \cong 0$ hence must be $0$ too. But if two out of three objects in a distinguished triangle are $0$ so must be the third and hence we find that $A_\infty\Dotimes_{(A, R)_\blacksquare} C \cong 0$ as well. Recall that $j_! j^\ast C \cong j_!A\Dotimes_{(A, R)_\blacksquare} C$ and that we have a distinguished triangle $A\to A_\infty \to A_\infty/A \to A[1]$ with $j_!A \cong (A_\infty/A)[-1]$. This gives a distinguished triangle
	$$j_! j^\ast C \to C\to A_\infty \Dotimes_{(A, R)_\blacksquare} C \to j_!j^\ast C[1]$$
	where again two of the three legs are $0$ so we find that also $C\cong 0$, \ie $C$ is acyclic as desired.
	
	Next, we will verify that $f^!R$ is a left bounded complex of discrete finitely generated $A$-modules. By \assumptionPreWord \ref{assumption: derived adjunctions are enriched} the adjunction $f_!\ladj f^!$ is $\condensedAb$-enriched which implies that
		$$\eRHom_R(f_!A, R)\cong \eRHom_A(A, f^!R)$$
	Thus $f^!R \cong \intRHom_R(f_!A, R)$ as $A$-modules as $\eRHom_R$ is the underlying (complex of) condensed abelian group(s) of $\intRHom_R$ and similarly the right-hand side is $\intRHom_A(A, f^!R) \cong f^!R$. Now $f_!A = (A_\infty/A)[-1]$ yields that
		$$f^!R \cong \intRHom_R((A_\infty/A)[-1], R) \cong \intRHom_R(A_\infty/A, R)[1].$$
	Now Clausen and Scholze state at the end of \cite[p. 57]{condenseddotpdf} that
		$$\intRHom_R(A_\infty/A, R)\cong A$$
	in case that the ring $R$ is $\Z$-finitely generated. The author was not able to verify that claim -- neither for a general $R$ nor in case of finite generation. We thus have to make it an \assumptionPreWord that $\intRHom_R(A_\infty/A, R)\cong A$. Using this one then immediately obtains that $f^!R \cong A[1]$ is a bounded complex of discrete finitely generated invertible $A$-modules.
		
	Using previous work it is easy to verify that $f_!$ preserves derived compactness. By lemma \ref{lemma: j lower shriek and compact generators} we have seen that $f_! =\scalarRestriction{\phi}\circ j_!$ preserves the derived compact generators of $\derived{A_\blacksquare}$. As $f_!$ commutes with arbitrary direct sums this is enough to see that $f_!$ preserves derived compactness in general.

	Thus, for the current case it remains to see that $f^!$ is a twist. Let $M\in\derived{R_\blacksquare}$. Using the component of the counit $f_!f^! R \to R$ and the projection formula we obtain a morphism
		$$f_!(f!R \Dotimes_{A_\blacksquare} (A_\blacksquare\Dotimes_{R_\blacksquare} M))\cong f!f^! R\Dotimes_{R_\blacksquare} M \to R\Dotimes_{R_\blacksquare} M \cong M.$$
	Under the adjunction $f_!\ladj f^!$ this morphism corresponds to a morphism
		$$f^!R\Dotimes_{A_\blacksquare}(A_\blacksquare\Dotimes_{R_\blacksquare} M) \to f^!M.$$
	If this is an isomorphism for every $M$ then we have shown that $f^!$ is indeed a twist. Now as both sides commute with arbitrary direct sums it is enough to show that this morphism is an isomorphism on generators $\free{R_\blacksquare}{S}$ for $S$ extremally disconnected. Now the left side is
		$$A[1]\Dotimes_{A_\blacksquare} \free{A_\blacksquare}{S} \cong \free{A_\blacksquare}{S}[1]$$
	as $f^!R \cong A[1]$ and $A_\blacksquare\Dotimes_{R_\blacksquare} \free{R_\blacksquare}{S}\cong \free{A_\blacksquare}{S}$. The author is unable to proceed from here. In \cite[p. 58]{condenseddotpdf} Clausen now continue by using Specker's theorem by writing $\free{R_\blacksquare}{S}\cong \prod_I R$ for some set $I$ -- this in not possible here since $R$ is not necessarily $\Z$-finitely generated. Then indeed $f^!\free{R_\blacksquare}{S} \cong \prod_I f^!R \cong \prod_I A[1] \cong \free{A_\blacksquare}{S}$ as $f^!$ was shown to commute with direct products. The author is not able to reproduce that $f^!\free{R_\blacksquare}{S} \cong \free{A_\blacksquare}{S}[1]$ in case that $R$ is not $\Z$-finitely generated. We thus have to make the \assumptionPreWord that $f^!\free{R_\blacksquare}{S} \cong \free{A_\blacksquare}{S}[1]$ for any extremally disconnected $S$ in case that $R$ is not $\Z$-finitely generated.
	
	\textbf{The case of a surjection $R\twoheadrightarrow A$:}\\
	If $R\to A$ is surjective then $\id_A\colon (A, R)_\blacksquare\to A_\blacksquare$ is an isomorphism by lemma \ref{lemma: analytic rings associated to finite morphisms}, hence $j_!$ is simply the identity and the projection formula easily follows.
	
	Now by assumption \ref{assumption: derived adjunctions are enriched} we find that the adjunction $f_!\ladj f^!$ is $\condensedAb$-enriched yielding that
		$$\eRHom_R(\scalarRestriction{\phi}A, R) \cong \eRHom_A(A, f^!R).$$
	Thus, $f^!R \cong \intRHom_R(A, R)$ as $A$-modules as $\eRHom_R$ is the underlying (complex of) condensed abelian group(s) of $\intRHom_R$ and similarly the right-hand side is $\intRHom_A(A, f^!R)\cong f^!R$. Thus, by replacing $A$ with a left resolution of discrete finitely generated projective $R$-modules we find that the underlying (complex of) $R$-module(s) of $f^!R$ is bounded to the left and consists of discrete finitely generated modules. The same is then true for the $A$-module $f^!R$ as well.
	
	Now assume that $f$ is of finite Tor-dimension. Then by \ref{lemma: agreement of dimensions} we can even choose thus resolution of $A$ to be bounded, yielding that $f^!R$ is bounded too. As scalar restriction $\scalarRestriction{\phi}$ preserves compactness by lemma \ref{lemma: scalar restriction preserves compactness for finite Tor-dimension} we find that $f_!=\scalarRestriction{\phi}\circ j_!$ preserves compactness as well since again, $j_!$ is trivial.
	
	Thus, for this case it remains to show that $f^!$ is a twist. As in the previous case where $A$ was a free $R$-algebra, the existence of the morphism
		$$f^!R\Dotimes_{A_\blacksquare}(A_\blacksquare\Dotimes_{R_\blacksquare} M) \to f^!M$$
	natural in $M\in \derived{R_\blacksquare}$ follows from the projection formula and the adjunction $f_!\ladj f^!$. The remaining argument analogous to that case as well. In particular, we also have to make the \assumptionPreWord that $f^!\free{R_\blacksquare}{S}\cong \free{A_\blacksquare}{S}[1]$ for any extremally disconnected $S$ here too. If $R$ is $\Z$-finitely generated then this is true.
	
	\textbf{The case of a complete intersection:}\\
	$f$ complete intersection:\\
	Now suppose that $R\to A$ is a even a complete intersection so that $A \cong R/\ideal{f_1,\dots, f_r}$ where $s=(f_1,\dots,f_r)$ is a regular sequence in $R$. Then by \cite[Cor. 4.5.5]{weibel-homological-algebra} the \emph{Koszul complex}
		$$K(s) = \mleft(0\to \exterior^r R^r \to \cdots \to \exterior^1 R \to \exterior^0 R^r \to 0\mright)$$
	(with $\exterior^0R^r$ in degree $0$) is an $R$-free left resolution of $A \cong R/\ideal{f_1,\dots, f_r}$. Since $R\to A$ is surjective we already have calculated that $f^R\cong \intRHom_R(A, R)$. It follows that
		$$f^!R \cong \intRHom_R(A, R) \cong \intRHom_R(K(s), R) \cong K(s)[-r]\cong A[-r]$$
	is an invertible $A$-module.
\end{proof}



\section{Affine duality rephrased geometrically}

First we want to observe that for an affine scheme the quasi-coherent modules embed well into solid modules. For this we first need that discrete modules are solid. We already have used this occasionally but for completeness we will give an argument here.
\begin{lemma}[Discrete modules are solid]
	\label{lemma: discrete modules are solid}
	
	Let $A$ be a discrete ring. If $M$ is a discrete $A$-module then the discrete condensed $\associated{A}$-module $\associated{M}$ is $A_\blacksquare$-complete. The natural inclusion
		$$\mod{A}\to\mod{A_\blacksquare}$$
	is a fully faithful and exact functor of abelian categories.
\end{lemma}
\begin{proof}
	Choose a presentation $A^J\to A^I\to M\to 0$ of $M$. Then $\associated A^J\to\associated A^I \to \associated M \to 0$ is a presentation of $\associated M$. Now $\associated A\cong A_\blacksquare[\ast]$ is $A_\blacksquare$-free, so $A_\blacksquare$-complete. Thus, both $\associated A^J$ and $\associated A^I$ are $A_\blacksquare$-complete and so is $M$ as a cokernel between $A_\blacksquare$-completes.
	
	As any discrete space is \CGWH we find that the inclusion is fully faithful by theorem \ref{theorem: some correspondences}. As limits and colimits of condensed $A$-modules are calculated locally it follows that the inclusion is exact.
\end{proof}

\begin{remark}
	Let $X=\Spec A$ be an affine scheme. In a geometric context we might call an object $\sheaf F\in \mod{A_\blacksquare}$ a \emph{solid condensed $\struct{X}$-module}. From remark \ref{remark: quasi-coherent sheaves of affine schemes} we know that we have an exact adjoint equivalence
		$$\mod{A} \xleftrightarrows{\Gamma(-, X)}{\associatedSheaf{(-)}} \qCoh{\struct{X}}$$
	asserting that the quasi-coherent $\struct{X}$-modules are precisely $A$-modules. By lemma \ref{lemma: discrete modules are solid} we thus obtain a functor
		$$\qCoh{X} \to \mod{A_\blacksquare}$$
	exhibiting quasi-coherent $\struct{X}$-modules as a full abelian subcategory of $\mod{A_\blacksquare}$ -- the category of \emph{solid condensed $\struct{X}$-modules}. Furthermore, if we denote by $\derived{\struct{X, \blacksquare}} \ldef \derived{A_\blacksquare}$ the derived category then we obtain that
		$$\derived{\qCoh{X}} \to \derived{\struct{X, \blacksquare}}$$
	\ie that the derived quasi-coherent $\struct{X}$-modules embed into $\derived{\struct{X, \blacksquare}}$.
\end{remark}

We are now close to stating the analogue of coherent duality \ref{theorem: coherent duality} in this extended setting. First, we need to investigate properness. In the affine case \emph{being proper} is a very strong condition. Any fiber of an affine morphism is again affine -- but fibers of a proper morphism should be \quote{complete}. Since an affine scheme is only complete if it is a finite union of points, one imagines that the proper morphism must already be finite.
\begin{lemma}[Proper morphisms of affine schemes are finite]
	If a morphism of affine, locally Noetherian schemes $\Spec A\to \Spec R$ is proper then it is finite, that is the $R$-module $A$ is finitely generated.
\end{lemma}
\begin{proof}
	By \cite[\href{https://stacks.math.columbia.edu/tag/02JJ}{Tag 02JJ}]{stacks-project} the morphism $R\to A$ is finite if and only if it is integral and of finite type. By definition, a proper morphism is of finite type. It remains to show that $R\to A$ is integral. By \cite[\href{https://mathoverflow.net/a/125793/546808}{this MathOverflow answer}]{proper-affine-integral} this follows as well.
\end{proof}


\begin{remark}
	The definition of $f^!$ as a right-adjoint to $f_!$ enriches by \assumptionPreWord \ref{assumption: derived adjunctions are enriched} yielding that for any $M\in D(A_\blacksquare)$ one has
		$$\eRHom_R(f_!M,R)\cong \eRHom_A(M, f^!R)$$
	inside $\derived{\condensedAb}$. As these are the underlying (complexes of) condensed abelian group(s) of $\intRHom_R$ and $\intRHom_A$ we find further that
		$$\intRHom_R(f_!M, R)\cong \RightD f_\ast\intRHom_A(M, f^!R)$$
	as $\RightD f_\ast \cong f_\ast$ as scalar restriction is exact.
	
	If $f$ is a complete intersection the complex $\omega_A\ldef f^!R$ is invertible. If $P$ is a finite projective $A$-module (or more generally a perfect complex of $A$-modules), then for $M=\RightD\intHom_{\associated{R}}(P, \omega_A)\cong P^\dual\otimes_A \omega_A$ this gives an isomorphism
	$$\RightD\intHom_{\associated{R}}(f_!(P^\dual\otimes_A \omega_A), R)\cong \RightD\intHom_{\associated{A}}(M, \omega_A) \cong P$$
	in $\derived{{\associated{R}}}$.
\end{remark}

\begin{theorem}[Solid coherent duality]
	\label{theorem: solid coherent duality}
	
	Let $f\colon X\to Y$ be a proper morphism of finite type between locally Noetherian affine schemes. Then $\RightD f_\ast\colon \derived{\struct{X,\blacksquare}} \to \derived{\struct{Y,\blacksquare}}$ agrees with the functor $f_!$ and hence admits a right adjoint $f^!\colon \derived{\struct{Y, \blacksquare}}\to \derived{\struct{X,\blacksquare}}$. Furthermore, there is an isomorphism
		$$\RightD f_\ast \intRHom_\struct{X}(\sheaf{F}, f^!\sheaf{G}) \cong \intRHom_{\struct{Y}}(\RightD f_\ast \sheaf F, \sheaf{G})$$
	of derived solid condensed $\struct{Y}$-modules natural in the derived solid condensed modules $\sheaf F\in\derived{\struct{X,\blacksquare}}$ and $\sheaf{G}\in\derived{\struct{Y,\blacksquare}}$.
\end{theorem}

\begin{remark}
	In the setting of theorem \ref{theorem: solid coherent duality} it is clear that $\RightD f_\ast = f_!$ preserves discrete objects -- that is the objects of $\qCoh{X}$. Let $X=\Spec A$ and $Y=\Spec R$. Now if $f$ is of finite Tor-dimension then theorem \ref{theorem: exceptional direct and inverse image functors, the affine case} implies that $f^!$ is a twist $f^!R \Dotimes_{A_\blacksquare}(A_\blacksquare\Dotimes_{R_\blacksquare} -)$ of (derived completed) scalar extension and that $f^!R$ is a bounded complex of discrete finitely generated $A$-modules. But this scalar extension is the unique colimit preserving extension of $\free{R_\blacksquare}{S}\mapsto \free{A_\blacksquare}{S}$ so in particular maps $R\cong\free{R_\blacksquare}{\ast}$ to $A\cong \free{A_\blacksquare}{\ast}$ and hence preserves discrete objects as well. Now finally since $\Dotimes_{A_\blacksquare}$ is the completion of the ordinary tensor product of $A$-modules we find that $f^!$ too preserves discrete objects. This implies that we have an adjunction
		$$\derived{\qCoh{Y}}\xleftrightarrows{\RightD f_\ast}{f^!}\derived{\qCoh{X}}$$
	and the corresponding isomorphism
		$$\RightD f_\ast \intRHom_\struct{X}(\sheaf{F}, f^!\sheaf{G}) \cong \intRHom_{\struct{Y}}(\RightD f_\ast \sheaf F, \sheaf{G})$$
	of $\struct{Y}$-modules natural in the sheaves $\sheaf F\in\derived{\qCoh{X}}$ and $\sheaf{G}\in\derived{\qCoh{Y}}$. As this adjunction even restricts to coherent sheaves (indeed $f^!$ is a twist by a complex of finitely generated $A$-modules so a twist by a coherent sheaf) we recover classic coherent duality for affine schemes as in theorem \ref{theorem: coherent duality} -- if $f$ is of finite Tor-dimension. This is often the case, \eg when $Y$ is regular (this includes the case where the base $Y$ is $\Spec K$ for some field) or when $f$ is a complete intersection (the Koszul complex gives at least one finite flat resolution of $A$).
\end{remark}
